\chapter{Computational Experiments}\label{chapter3}

This chapter contains the computational experiments for the algorithms proposed in comparison with the current state of the art algorithms. To be precise, I will compare the ACO and genetic algorithm proposed with the branch-and-cut algorithms CEC and CUT from Marzo et al.~\cite{marzo2} and the HLIPP heuristic from Marzo and Ribeiro~\cite{marzo}. All algorithms run on the same hardware with the same amount of time limit and with the same resources. The exact algorithms failed to produce any significant results on larger graphs to be included in the comparison. Their results are still present in the papers referenced throughout this thesis as most of the instances are shared.

\section{Test Instances}\label{sec:instances}

The test instances used are real world network instances from Marzo and Ribeiro~\cite{marzo} which are the standard in the literature around this problem. They range from social networks with 33 networks to the yeast protein interaction network which has 2361 vertices and 6646 edges. We know the optimal values for most of these instances using exact algorithms or from prior work. The yeast instance has no known optimal value as no exact algorithm has solved it within practical time limits. I've also generated a synthetic instance with 5000 vertices and over 1.300.000 edges with a known lower bound on the optimal solution using the generator described in the next section. This instance is meant as a very hard instance to solve for any algorithm and is used to compare the algorithms on dense graphs.

\subsection{Instance generator}\label{sec:generator}

For real-world large instances which are hard to solve, the optimal solution is unknown so there is no way to measure how close each algorithm gets to finding the optimal solution. To help with this, I've built an instance generator capable of generating hard to solve instances with a known lower bound for the optimal path. While there exists a chance for an induced path to exist that is longer than the known lower bound, the difference would be at most a few vertices. The algorithm consists of a planted path of a fixed size then adding new vertices and edges in such a way to complicate any algorithm trying to find it. At the end, the algorithm does a double check to check that the path is still induced.

The generator has four steps:

\begin{enumerate}
    \item \textbf{Planted path.} A simple path of length $L$ is created using vertices $\{0, 1, \ldots, L-1\}$ with edges connecting consecutive vertices.

    \item \textbf{Noise vertices.} Create $N - L$ vertices as noise then we connect them to a percentage of the vertices in the real path (depending on the difficulty setting, from 20-35\%). Since each noise vertex touches many path vertices, including it in an induced path forces the exclusion of more path vertices than it contributes which guarantees that it is always a net loss.

    \item \textbf{Clique structure.} Noise vertices are organized into small cliques of size 3, which limits any induced path through the noise subgraph to at most 2 vertices per clique. We also add edges between cliques with a low probability (10\%) to prevent the noise structure from being easily separable.

    \item \textbf{Trap structures.} Add short paths branching off near the endpoints of the planted path. These paths are traps which look like promising extensions but lead into the clique-structured noise region where further extension is quickly blocked.
\end{enumerate}

The generator supports three difficulty levels (EASY, MEDIUM, HARD) which control the noise-to-path connection density, the clique size, and the trap parameters. The highest difficulty (HARD) connects each noise vertex connects to a smaller fraction of path vertices, and more and longer traps are added.

The instance used in the experiments, \textit{generated-hard-5000}, was created with HARD parameters and it contains 5000 vertices, a planted induced path of length 1600, the noise-to-path density of 22\%, 10\% inter-clique edges probability and 6 traps of length 500. The resulting graph has 1,310,158 edges and an average degree of approximately 524 which makes it far denser and harder to solve than any of the real-world instances.


\section{Computational Environment}\label{sec:environment}

The algorithms were implemented in Java 21 using the Graph4J for graph storage and manipulation. The experiments ran on a laptop with a 5800H 3.2 GHz CPU and 16 GB of RAM running on Windows 11 Pro 64-bit. Gurobi 13.0.1 was used to run any integer formulation presented in the experiments.

The branch and cut algorithms (CEC and CUT) and the HLIPP heuristics were implemented based on the instructions and pseudocode in in~\cite{marzo2} and~\cite{marzo} respectively, I've double checked the results I've obtained against those reported in the original publications to verify correctness.

Each algorithm has been given a time limit of 15 minutes per instance. The algorithms that have a random factor involved (GA and ACO) were run 5 times per instance while the deterministic algorithms (CEC, CUT and HLIPP) were run only once. For the genetic algorithm and ACO, I've reported the min, average, median and maximum values for a correct comparison. All algorithms were ran on the same computational environment with the same resources.

\section{Results}\label{sec:results}

The table below presents the results obtained from running all the algorithms on the 25 test instances. The first four columns represents the details about the instance: the name, number of vertices (\(|V|\)), the number of edges (\(|E|\) and the optimum induced path length (\(OPT\)) if it is known. The rest of the columns are the values obtained by the current state of the art algorithms and the maximum values found by the proposed algorithms. If the value is the optimal value then they are shown in \textbf{bold}. A dash (--) indicates that the algorithm failed to produce a solution within the time limit.

\linespread{1.0}
\begin{table}[H]
\centering
\caption{Results and running times of all algorithms on all test instances. GA and ACO results are averages over 5 runs, rounded to the nearest integer; times are the average running time. Bold values match the proven optimal.}
\label{tab:results}
\resizebox{\textwidth}{!}{%
\setlength{\tabcolsep}{4pt}
\begin{tabular}{lccr|rr|rr|rr|rr|rr}
\toprule
\multirow{2}{*}{Instance} & \multirow{2}{*}{$|V|$} & \multirow{2}{*}{$|E|$} & \multirow{2}{*}{OPT} & \multicolumn{4}{c|}{Marzo et al.~\cite{marzo2}} & \multicolumn{2}{c|}{Marzo~\cite{marzo}} & \multicolumn{4}{c}{Proposed} \\
\cmidrule(lr){5-8} \cmidrule(lr){9-10} \cmidrule(lr){11-14}
& & & & CEC & Time & CUT & Time & HLIPP & Time & GA & Time & ACO & Time \\
\midrule
high-tech        & 33   & 91       & 13   & \textbf{13} & 0.2   & \textbf{13} & 0.1   & \textbf{13} & 0.1   & \textbf{13} & 11.6  & \textbf{13} & 0.3  \\
karate           & 34   & 78       & 9    & \textbf{9}  & 0.1   & \textbf{9}  & 0.1   & \textbf{9}  & 0.0   & \textbf{9}  & 10.2  & \textbf{9}  & 0.2  \\
mexican          & 35   & 117      & 16   & \textbf{16} & 0.2   & \textbf{16} & 0.2   & \textbf{16} & 0.1   & \textbf{16} & 13.9  & \textbf{16} & 0.2  \\
sawmill          & 36   & 62       & 18   & \textbf{18} & 0.0   & \textbf{18} & 0.0   & \textbf{18} & 0.0   & \textbf{18} & 12.6  & \textbf{18} & 0.1  \\
tailorS1         & 39   & 158      & 13   & \textbf{13} & 0.2   & \textbf{13} & 0.2   & \textbf{13} & 0.1   & \textbf{13} & 15.6  & \textbf{13} & 0.3  \\
chesapeake       & 39   & 170      & 16   & \textbf{16} & 0.2   & \textbf{16} & 0.2   & \textbf{16} & 0.1   & \textbf{16} & 15.4  & \textbf{16} & 0.3  \\
tailorS2         & 39   & 223      & 15   & \textbf{15} & 0.6   & \textbf{15} & 0.6   & \textbf{15} & 0.2   & \textbf{15} & 16.0  & \textbf{15} & 0.3  \\
romeo and juliet & 41   & 120      & 9    & \textbf{9}  & 0.0   & \textbf{9}  & 0.0   & \textbf{9}  & 0.0   & \textbf{9}  & 11.0  & \textbf{9}  & 0.2  \\
die hard         & 47   & 237      & 10   & \textbf{10} & 0.1   & \textbf{10} & 0.1   & \textbf{10} & 0.0   & \textbf{10} & 14.3  & \textbf{10} & 0.4  \\
attiro           & 59   & 128      & 31   & \textbf{31} & 0.2   & \textbf{31} & 0.1   & 30          & 1.6   & \textbf{31} & 26.9  & \textbf{31} & 0.3  \\
krebs            & 62   & 153      & 17   & \textbf{17} & 0.2   & \textbf{17} & 0.2   & \textbf{17} & 0.2   & \textbf{17} & 23.2  & \textbf{17} & 0.3  \\
dolphins         & 62   & 159      & 24   & \textbf{24} & 0.2   & \textbf{24} & 0.2   & 23          & 1.4   & \textbf{24} & 28.4  & \textbf{24} & 0.3  \\
prison           & 67   & 142      & 36   & \textbf{36} & 0.0   & \textbf{36} & 0.0   & \textbf{36} & 2.6   & \textbf{36} & 32.3  & \textbf{36} & 0.3  \\
huck             & 69   & 297      & 9    & \textbf{9}  & 0.1   & \textbf{9}  & 0.4   & \textbf{9}  & 0.1   & \textbf{9}  & 20.2  & \textbf{9}  & 0.8  \\
sanjuansur       & 75   & 144      & 38   & \textbf{38} & 0.2   & \textbf{38} & 0.1   & 36          & 2.6   & \textbf{38} & 39.6  & \textbf{38} & 0.3  \\
jean             & 77   & 254      & 11   & \textbf{11} & 0.1   & \textbf{11} & 0.2   & \textbf{11} & 0.3   & \textbf{11} & 26.2  & \textbf{11} & 0.6  \\
david            & 87   & 406      & 19   & \textbf{19} & 0.1   & \textbf{19} & 0.1   & \textbf{19} & 2.0   & \textbf{19} & 36.3  & \textbf{19} & 1.2  \\
ieeebus          & 118  & 179      & 47   & \textbf{47} & 0.3   & \textbf{47} & 0.2   & \textbf{47} & 8.2   & \textbf{47} & 62.0  & \textbf{47} & 0.3  \\
sfi              & 118  & 200      & 13   & \textbf{13} & 0.0   & \textbf{13} & 0.0   & \textbf{13} & 0.1   & \textbf{13} & 38.1  & \textbf{13} & 0.8  \\
anna             & 138  & 493      & 20   & \textbf{20} & 0.2   & \textbf{20} & 0.2   & \textbf{20} & 2.5   & \textbf{20} & 57.8  & \textbf{20} & 2.2  \\
\midrule
usair            & 332  & 2\,126   & 46   & \textbf{46} & 123.7 & \textbf{46} & 104.2 & 38          & 50.8  & \textbf{46} & 242.9 & \textbf{46} & 16.2 \\
494bus           & 494  & 586      & 142  & \textbf{142}& 16.6  & \textbf{142}& 88.5  & 109         & 220.9 & \textbf{142}& 561.7 & \textbf{142}& 2.5  \\
662bus           & 662  & 906      & 305  & \textbf{305}& 127.0 & \textbf{305}& 811.0 & 237         & 265.2 & 277         & 900.0 & 298         & 900.0  \\
yeast            & 2\,361 & 6\,646 & ?    & 142         & 900.0 & 184         & 900.0 & 204         & 900.0 & 300         & 900.0 & 387         & 900.0\\
gen-hard-5000    & 5\,000 & 1\,310\,158 & $\geq$1600\textsuperscript{$\dagger$} & 51 & 900.0 & 51 & 900.0 & 313 & 900.0 & 387 & 900.0 & 336 & 900.0 \\
\bottomrule
\end{tabular}%
}
\vspace{0.2cm}

\raggedright
\footnotesize
\textsuperscript{$\dagger$} Lower bound from planted induced path (see Section~\ref{sec:generator}). Times are in seconds.
\end{table}
\linespread{1.5}

\subsection{Small and medium instances}

On all instances up to anna (138 vertices, 493 edges), every algorithm finds the optimal solution. CEC and CUT solve these in under a second in most cases. HLIPP also finds the optimal on most instances but misses it on three: attiro (30 vs 31), dolphins (23 vs 24) and sanjuansur (36 vs 38). The genetic algorithm and the ant colony optimization algorithm find the optimal on all small and medium instances consistently across all 5 runs.

\subsection{Large real-world instances}

The differences between algorithms become apparent on the larger instances. Table~\ref{tab:detailed} shows the minimum, median, and maximum results over 5 runs for the stochastic algorithms, along with the average running time.

\linespread{1.0}
\begin{table}[H]
\centering
\caption{Detailed results on large instances. CEC, CUT, and HLIPP are single-run deterministic results. GA and ACO show min/average/max over 5 runs. All times are in seconds.}
\label{tab:detailed}
\hspace*{-1.05cm}
\setlength{\tabcolsep}{3pt}
\footnotesize
\begin{tabular}{lc|rr|rr|rr|rrrr|rrrr}
\toprule
& & \multicolumn{4}{c|}{Marzo et al.~\cite{marzo2}} & \multicolumn{2}{c|}{Marzo~\cite{marzo}} & \multicolumn{4}{c|}{GA} & \multicolumn{4}{c}{ACO} \\
Instance & OPT & CEC & Time & CUT & Time & HLIPP & Time & Min & Avg & Max & Time & Min & Avg & Max & Time \\
\midrule
usair         & 46  & \textbf{46} & 123.7 & \textbf{46} & 104.2 & 38 & 50.8 & 45  & 45.8  & \textbf{46}  & 242.9  & \textbf{46}  & \textbf{46.0}  & \textbf{46}  & 16.2   \\
494bus        & 142 & \textbf{142} & 16.6 & \textbf{142} & 88.5 & 109 & 220.9 & \textbf{142} & \textbf{142.0} & \textbf{142} & 561.7  & \textbf{142} & \textbf{142.0} & \textbf{142} & 2.5    \\
662bus        & 305 & \textbf{305} & 127.0 & \textbf{305} & 811.0 & 237 & 265.2 & 268 & 277.0 & 283 & 900.0  & 287 & 297.6 & 302 & 900.0    \\
yeast         & ?   & 142 & 900.0 & 184 & 900.0 & 204 & 900.0 & 279 & 300.2 & 328 & 900.0  & 384 & 387.4 & 395 & 900.0  \\
gen-hard-5000 & $\geq$1600 & 51 & 900.0 & 51 & 900.0 & 313 & 900.0 & 375 & 386.8 & 395 & 900.0 & 334 & 336.2 & 339 & 900.0  \\
\bottomrule
\end{tabular}
\vspace{0.2cm}

\raggedright
All algorithms on these instances reached the 15-minute time limit.
\end{table}
\linespread{1.5}

On \textbf{usair} (332 vertices, 2126 edges), all algorithms find the optimal solutipn except the HLIPP algorithm. The fastest among these is the ACO which took 16.2 seconds. CUT algorithm took 104 seconds and CEC 123 secondes respectively. The genetic algorithm also managed to obtain the optimum value 4 times out of 5 in 242 seconds.

On \textbf{494bus} (494 vertices, 586 edges), CEC gets the optimal value of 142 in 16.6 seconds and CUT in 88.5 seconds. HLIPP finds only 109. Both GA and ACO find the optimal on all 5 runs, with ACO converging in 2.5 seconds on average compared to 561.7 for GA.

On \textbf{662bus} (662 vertices, 906 edges), CEC and CUT manage to find the optimal value of 305 which was previously unknown in the literature. CEC finishes in 127 seconds while CUT needs 811 seconds. HLIPP reaches only 237. The best value ACO finds is 302 (with an avg of 297.6) missing the optimal by only 3 vertices. The best value found by the genetic algorithm is 283 further from the optimal. This instance is the last one where CEC and CUT manage to get good results. Till to this point, the CEC and CUT are the best algorithms but adding further vertices and edges start to severely impact the performance of the branch and cut algorithms.

On \textbf{yeast} (2361 vertices, 6646 edges), no exact algorithm finishes within the time limit: CEC finds only 142 and CUT finds 184 before timing out. HLIPP reaches only 204 before timing out. The best value found by the Ant Colony Optimization algorithm is 395 with an average of 387.4, which is the best known result for this instance. The best found by the GA is 328 with an avg of 300.2, well below ACO which performs the best for this instance.

\subsection{Generated instance}

The generated instance is the hardest instance of them all with 5000 vertices and 1,310,158 edges with an average degree of 524. This is much denser than any of the real world instances and contains many clique structures and a lot of noise. The search space is large but very few of its paths are long.

The induced path planted for this instance is of length 1600. None of the algorithms even came close to this value which shows just how difficult such an instance really is. The longest path found by both the CEC and CUT algorithms is of length 51 before timing out. At this stage, any exact algorithm, even those using integer formulations are no longer feasible. The HLIPP algorithm reaches only 313 before timing out. The genetic algorithm managed to outperform the ACO algorithm obtaining the maximum value of 395 compared to 339 of the Ant Colony Optimization. This shows that each metaheuristic is best suited for a type of graphs, in this case, the genetic algorithm is better fitted for denser graphs.


\subsection{Algorithm comparison}\label{sec:comparison}

\paragraph{Exact algorithms (CEC and CUT).} The CEC and CUT algorithms both perform exceptionally well on all instances up to 662bus. The CEC implementation seems to be better as it takes less time. For smaller to medium instances, these two algorithms are the best performing. For larger instances though these algorithms failed to find good results.

\paragraph{HLIPP heuristic.} HLIPP finds the optimal for smaller instances but fails to provide good results for medium to large instances. The deterministic path enumeration becomes inefficient on large or dense graphs because the combinatorial explosion prevents it from reaching good regions of the search space.

\paragraph{ACO.} The Ant Colony Optimization is the best algorithm found in the experiments for sparse graphs such as the real-world instances used. The direction-based pheromone approach which also takes into account the length of the paths yielded good results. It also efficiently uses all resources by running in parallel as each ant uses one thread of the compute environment. The ACO managed to get the optimum value on all small and medium instances in a short amount of time. For larger instances, ACO managed to stay competitive on \textbf{662bus} obtaining 302, 3 vertices away from the optimu mvalue of 305. The algorithm managed to find the best value found so far for \texbf{yeast} instance obtaining 395 which is much higher than the 328 found by the genetic algorithm or the 184 and 142 values found by the CUT and CEC algorithms respectively.

\paragraph{GA.} The genetic algorithm proposed for this problem managed to find the optimal solution for all smaller to medium instances in an average amount of time. For larger instances that are sparse it performed less optimal than ACO as it obtained worse results in the same amount of time. The real value of the genetic algorithm lies in dense graphs, on the hardest instance with 5000 vertices and over 1,310,158 edges, the genetic algorithm managed to outperform all the other algorithms. The longest induced path found by the GA on this instance was 395 compared to the ACO's 339 and the 51 found by CEC and CUT meaning that the genetic algorithm is better on dense graphs.
